% lemmon-fitch.tex
%
% First draft. Compile with pdflatex. Packages: the AMS bundle, geometry,
% array, and listings -- all part of a standard TeX distribution. The proof
% displays are drawn with a small macro defined below rather than with
% fitch.sty, to keep the source portable; if you would rather use Peter
% Selinger's fitch.sty for the final version, the displays are all in one
% place and easy to convert.
%
% The development now has two machine-readable faces.  Representative Haskell
% excerpts explain the algorithms, while representative Isabelle/HOL excerpts
% record the definitions and theorems used to settle the mathematical claims.

\documentclass[11pt]{article}

\usepackage[margin=1.15in]{geometry}
\usepackage{amsmath,amssymb,amsthm}
\usepackage{array}
\usepackage{listings}

\lstset{
  language=Haskell,
  basicstyle=\small\ttfamily,
  keywordstyle=\bfseries,
  commentstyle=\itshape,
  columns=fullflexible,
  showstringspaces=false,
  xleftmargin=1.5em,
  aboveskip=1.2em,
  belowskip=1.2em,
  morekeywords={data,type,where,case,of,let,in,pure}
}

\theoremstyle{plain}
\newtheorem{theorem}{Theorem}
\newtheorem{proposition}[theorem]{Proposition}
\newtheorem{corollary}[theorem]{Corollary}
\newtheorem{lemma}[theorem]{Lemma}
\theoremstyle{definition}
\newtheorem{definition}[theorem]{Definition}
\newtheorem{example}[theorem]{Example}
\newtheorem{problem}[theorem]{Problem}
\theoremstyle{remark}
\newtheorem{remark}[theorem]{Remark}

% A scope bar for Fitch displays.
\newcommand{\fbar}{\rule[-0.7ex]{0.4pt}{2.6ex}\hspace{0.7em}}
\newcommand{\dep}[1]{\textup{#1}}

\newenvironment{lemmonproof}
  {\begin{center}\begin{tabular}{@{}r@{\hspace{1.5em}}r@{\hspace{0.6em}}l@{\hspace{2.5em}}l@{}}}
  {\end{tabular}\end{center}}

\newenvironment{fitchproof}
  {\begin{center}\begin{tabular}{@{}r@{\hspace{1em}}l@{\hspace{2.5em}}l@{}}}
  {\end{tabular}\end{center}}

\title{Dependency and Scope:\\
       On the Translation Between Lemmon and Fitch Proofs}
\author{Hans Halvorson}
\date{Revised draft with Isabelle/HOL verification}

\begin{document}
\maketitle

\begin{abstract}
Lemmon-style and Fitch-style natural deduction are usually treated as
notational variants of one another. They are, but the sense in which they are
is more delicate than that description suggests, and the delicacy is
instructive. A Lemmon line carries an explicit and exact record of the
assumptions it rests on; a Fitch line carries no such record, and its
dependencies are bounded, but not determined, by the subproofs enclosing it.
We define a canonical map $\delta$ from Fitch syntax to Lemmon syntax and
prove that every correct Fitch proof is sent to a correct Lemmon proof.
The converse implication, as originally formulated, is false: because
$\delta$ forgets excess scope, it can turn an eigenvariable-invalid Fitch
object into a correct Lemmon proof. A separate semantic audit found that the
Fitch turnstile, as first formalised, admitted a proof ending in an
undischarged top-level subproof, and so derived an arbitrary formula from no
premises. The closed-root condition this shows to be necessary is now part of
the definition, and the Fitch turnstile is proved semantically sound. These
machine-found boundary cases are stated explicitly rather than hidden by a
round-trip test. Translation in the other direction is not
\emph{positional}: there are Lemmon proofs whose lines cannot be carried over
one for one, in order, to any Fitch proof. The obstruction is structural.
Each notation presents a derivation---which is a tree---as a directed acyclic
graph, by writing each line once and citing it thereafter; but where Lemmon
permits a line to be cited by any later line whatever, Fitch permits citation
only along the paths its nesting of subproofs allows.

For the compact reconstructed calculus we give a construction translating any
correct Lemmon proof by way of its unfolding. Fresh-name repair handles the
eigenvariable conditions for both universal introduction and existential
elimination. In Isabelle/HOL we prove that this repaired construction
terminates and yields a correct Fitch proof with the same conclusion. We also
answer two economy questions negatively for explicit citation-preserving
notions of image. A
five-line Lemmon proof can discharge two different assumptions at the same
conclusion line; no permutation, no duplication-free layout, and not even the
addition of arbitrary auxiliary target lines can represent those source lines
with their citation-bearing rules unchanged in a well-formed Fitch proof.
Laminar dependency regions do not suffice when the original numbering and
order are retained; this does not rule out repair by permutation.
The Isabelle definitions generate
a Haskell and SML kernel, so the executable checker and the proved equations
have a common source. The proof checker supplied in
\verb|reference/lemmon-checker-main/src| is authoritative for the rule language
of \emph{How Logic Works}. Isabelle therefore contains a separate, exact
\verb|hl_formula| and \verb|hl_justification| layer matching its datatypes,
including Boolean truth constants, equality as a predicate, modus tollens,
excluded middle, propositional consequence, and quantifier negation. The
formal mirror also totalises Boolean and biconditional cases at which the
handwritten program was partial. Its general Lemmon and Fitch semantic
soundness results and its checked-output guarantees do not yet supply the
general unfolding-to-Fitch correctness theorem proved for the compact calculus.
\end{abstract}

\section{Introduction}

Anyone who has taught both notations knows that a Lemmon proof and a Fitch
proof of the same argument look like transcriptions of one another. In the
Lemmon style, each line is annotated with the numbers of the assumptions it
depends on, and a rule such as conditional proof discharges an assumption by
removing its number from the annotation. In the Fitch style, the same
information is carried geometrically: an assumption opens a subproof, drawn as
a vertical scope line, and discharging it closes the subproof. The two
conventions are so evidently doing the same work that the translation between
them is generally left as an exercise, or not mentioned.

It is not quite an exercise. The purpose of this paper is to say exactly what
the translation is, where it is easy, where it is not, and what the difficulty
reveals about the relationship between the two notations.

The central observation is simple enough to state here. In a Lemmon proof, the
dependency set attached to a line is \emph{exact}: it is the set of
assumptions the line actually rests on. In a Fitch proof, the corresponding
information is the set of assumptions of the enclosing subproofs, and this is
merely an \emph{upper bound}. A Fitch line may perfectly well sit inside a
subproof whose assumption it never uses. Every asymmetry we discuss below
follows from this one.

A second observation, less obvious, does most of the work in
Section~\ref{sec:obstructions}. In a Lemmon proof, the textual position of a
line carries no semantic content. Lines may be permuted freely, subject only
to the requirement that a line be written after the lines it cites. In a Fitch
proof, position is semantic: where a line sits determines what may see it, and
a line inside a subproof becomes unavailable once that subproof closes. The
translation from Lemmon to Fitch must therefore invent positional information
that the source does not contain, and it is possible for the source to place
demands that no such invention can satisfy.

A third observation concerns the eigenvariable rules. Universal introduction
and existential elimination impose freshness conditions not merely on the
formulas immediately cited but on assumptions. Lemmon can state those
conditions exactly. Fitch has no record of what a line rests on, only of what
encloses it, and must state them about scope instead. Since scope may properly
include the dependency set, the two conditions come apart, and a naive
translation of a correct Lemmon proof may break a Fitch freshness condition.
Section~\ref{sec:trees} exhibits one instance and gives the verified repair.

The structural arguments were first mechanised with a compact internal rule
language. The proof checker in the supplied Haskell \verb|src| tree is the
authoritative roster intended for \emph{How Logic Works}
\cite{halvorson2020}. The current Isabelle development now represents that
roster directly, alongside the internal layer used by the established
translation proofs. The compact roster differs in rules as well as syntax:
the supplied system includes modus tollens, propositional consequence,
excluded middle and quantifier negation, and its biconditional elimination has
a different citation pattern. There is no proved bridge transferring the
compact construction theorem to that exact executable language.
The general translation and image-nonexistence results below refer to the
compact calculus unless an exact \verb|HL_| predicate is specified.
Section~\ref{sec:machine} explains the correspondence and the defects that
total formal functions made visible. The separate comparison document
\verb|ORIGINAL_PAPER_COMPARISON.md| records every numbered item of the original
fourteen-page draft; the present file is a revised account, not that baseline.
Section~\ref{sec:systems} fixes notation.
Section~\ref{sec:fitch-to-lemmon} treats the easy direction and
draws from it the redundancy of the dependency column.
Section~\ref{sec:obstructions} exhibits the two obstructions to a positional
translation in the other direction. Section~\ref{sec:trees} introduces the
graph-theoretic vocabulary needed to say what the obstructions have in common,
locates the difference between the notations in how each constrains the
sharing of lines, gives the construction that translates by unfolding, and
repairs it. Section~\ref{sec:open} reports the Isabelle proofs and the negative
answers to the two economy questions. Section~\ref{sec:machine} explains the
formalisation and generated code, and
Section~\ref{sec:equivalence} draws the moral.

\section{The two systems}\label{sec:systems}

\begin{definition}
A \emph{Lemmon proof} is a finite sequence of lines. A line is a quadruple
$\langle \Gamma, n, \varphi, j\rangle$ where $n$ is the line's number,
$\varphi$ its formula, $j$ its justification, and $\Gamma$ a set of line
numbers---the \emph{dependency set}. A justification is either
$\mathrm{A}$ (assumption) or a rule together with the numbers of the lines it
cites, all of which must be smaller than $n$.
\end{definition}

The rules fall into two groups. Most of them---modus ponens, conjunction
introduction, universal elimination---simply pool the dependency sets of the
lines they cite. Four of them discharge: conditional proof and reductio each
cite an assumption line together with a line derived from it, and disjunction
and existential elimination each cite one or two such pairs. For these, the
discharged assumption is removed from the pooled set. Thus in

\begin{lemmonproof}
\dep{1}   & (1) & $P$              & A \\
\dep{2}   & (2) & $Q$              & A \\
\dep{1,2} & (3) & $P \wedge Q$     & 1,2 $\wedge$I \\
\dep{1}   & (4) & $Q \to (P \wedge Q)$ & 2,3 CP \\
\end{lemmonproof}

\noindent line (4) discharges assumption (2), which accordingly leaves the
dependency set.

In the implementation a line is a record and a justification an enumeration,
one constructor per rule. The constructors of the discharging rules are worth
noticing: each names an assumption line together with a line derived from it.

\begin{lstlisting}
data ProofLine = ProofLine
  { lineNumber    :: Int
  , formula       :: PredFormula
  , justification :: Justification
  , references    :: Set Int          -- the dependency set
  }

data Justification
  = Assumption
  | MP Int Int                        -- modus ponens, two lines
  | AndIntro Int Int
  | CP Int Int                        -- assumption line, conclusion line
  | RAA Int Int
  | OrElim Int Int Int Int Int        -- disjunction, then two such pairs
  | ExistsElim Int Int Int
  | ForallIntro Int
  -- ... fifteen further rules
\end{lstlisting}

\begin{definition}\label{def:fitch}
A \emph{Fitch proof} is a finite sequence of \emph{items}, where an item is
either a line---a number, a formula and a justification---or a
\emph{subproof}, consisting of an assumption line together with a further
sequence of items. The \emph{assumption scope} of a line consists of the outer
premises together with the assumptions of the subproofs containing it. A line
may cite another only if the cited line lies
within its scope, that is, only if every subproof containing the cited line
also contains the citing line.
\end{definition}

The corresponding type is mutually recursive, an item being either a line or
a subproof containing further items. A rule that discharges cites a subproof,
which is written as the pair of its first and last line---the \verb|SubRef|
below---and this is exactly the pair that the Lemmon rule already names.

\begin{lstlisting}
type SubRef = (Int, Int)

data FitchItem
  = FLine Int PredFormula FitchRule
  | FSub Subproof

data Subproof = Subproof
  { subAssumeLine :: Int
  , subAssumeForm :: PredFormula
  , subBody       :: [FitchItem]
  }

data FitchRule
  = FPremise                          -- never discharged: outermost level
  | FAssume                           -- opens a subproof
  | FMP Int Int
  | FCP SubRef                        -- cites the subproof, not two lines
  | FOrE Int SubRef SubRef
  -- ... and so on
\end{lstlisting}

The same proof in Fitch:

\begin{fitchproof}
1 &                 $P$                  & Premise \\
2 & \fbar           $Q$                  & Assume \\
3 & \fbar           $P \wedge Q$         & $\wedge$I 1,2 \\
4 &                 $Q \to (P \wedge Q)$ & CP 2--3 \\
\end{fitchproof}

Two conventions deserve comment. First, an assumption never discharged is
written in Fitch as a premise at the outermost level, so the choice between
``premise'' and ``assumption'' is determined by whether the line is discharged
later; nothing in the Lemmon proof marks the difference locally. Second, we
take universal introduction to carry a side condition on the name generalised
upon, rather than to open a subproof with a flagged constant. Both conventions
are found in the literature; the first agrees with the condition already
imposed in the Lemmon system, and choosing it lets universal introduction
translate one line to one line.

\section{From Fitch to Lemmon}\label{sec:fitch-to-lemmon}

This direction is unproblematic, and it is worth seeing why, because the
reason is exactly what fails in the other direction.

\begin{definition}
Let $F$ be a Fitch proof. Define $\delta(F)$, the Lemmon proof induced by $F$,
by flattening $F$ into a sequence of lines in order and assigning to each line
$n$ the set $\Gamma(n)$ given by: $\Gamma(n) = \{n\}$ if $n$ is an assumption
or premise; $\Gamma(n) = \bigcup_{m} \Gamma(m)$ for $m$ ranging over the lines
cited, if the rule at $n$ discharges nothing; and for a discharging rule,
the same union with the discharged assumptions removed.
\end{definition}

The whole of $\delta$'s arithmetic is the following function, which computes
one line's dependency set given those of the lines before it. The
discharging cases subtract; every other rule takes the union of what it cites.

\begin{lstlisting}
depsOf :: Map Int (Set Int) -> FitchRule -> Int -> Set Int
depsOf deps r self =
  case r of
    FPremise               -> singleton self
    FAssume                -> singleton self
    FCP  (a, c)            -> delete a (look c)
    FRAA (a, c)            -> delete a (look c)
    FOrE d (a1,c1) (a2,c2) ->
      unions [ look d, delete a1 (look c1), delete a2 (look c2) ]
    FExistsE m (a, c)      -> look m `union` delete a (look c)
    _                      -> unions (map look (citedLines (toLemmonRule r)))
  where
    look n = findWithDefault empty n deps
\end{lstlisting}

\begin{theorem}\label{thm:total}
$\delta$ is a total function on Fitch syntax. If $F$ is a correct Fitch proof,
then $\delta(F)$ is a correct Lemmon proof.
\end{theorem}

\begin{proof}
The definition of $\Gamma$ is a recursion on the position of the line, which
terminates because a line may cite only earlier lines. That the resulting
annotation satisfies the Lemmon conditions on each rule is checked rule by
rule; the discharging rules are the only ones where anything happens, and for
those the Fitch scope condition guarantees that the discharged assumption is
available to be removed.  Isabelle proves the implication as
\verb|theorem_4_forward|.
\end{proof}

\begin{remark}\label{rem:theorem4-converse}
The converse formerly asserted here is false for the present formal definition
of correctness. The six-line Fitch object displayed in the proof of
Theorem~\ref{thm:refuted} violates universal introduction because $F(a)$ is in
scope at line (4). Its $\delta$-image is nevertheless the correct Lemmon proof
displayed immediately above it: recomputation gives line (4) dependency set
$\{1\}$, whose assumption does not contain $a$. Thus correctness of
$\delta(F)$ cannot imply correctness of $F$. This does not impugn the forward
translation; it shows that a map which forgets excess scope cannot reflect a
scope-sensitive eigenvariable condition.
\end{remark}

\begin{proposition}\label{prop:minimal}
For each line $n$ of $F$, the set $\Gamma(n)$ assigned by $\delta$ is
contained in the scope of $n$, and the containment may be proper.
\end{proposition}

The containment may be proper precisely when a line sits inside a subproof
whose assumption it does not use. This has a consequence worth recording.

\begin{corollary}\label{cor:notinjective}
$\delta$ is not injective. Two Fitch proofs that differ only in the placement
of a line within subproofs it does not depend on have the same image.
\end{corollary}

The round trip need not be the identity. The formal example
\verb|roundtrip_normalises| moves an unused line out of a wider subproof.
This example and noninjectivity do not establish a general normalization
theorem or a minimal-width property of the resulting subproofs.

The next observation is elementary but, we think, has not been stated. Because
the Lemmon rules determine the dependency set of each line exactly---a system
that permitted a line to cite more assumptions than it uses would not satisfy
Proposition~\ref{prop:minimal}---the dependency column is not independent
data.

\begin{proposition}\label{prop:redundant}
Let $L$ be a Lemmon proof and let $L^{-}$ be the result of deleting its
dependency column. Then $L$ is recoverable from $L^{-}$.
\end{proposition}

\begin{proof}
The recursion in the definition of $\delta$ makes no use of the dependency
column of the source; it uses only the justifications. Applying it to $L^{-}$
returns $L$.
\end{proof}

This argument assumes exact dependency arithmetic. The faithful Haskell
checker does not impose it on excluded middle, and even
\verb|hlVerifiedCorrect| permits that rule to carry extra dependencies on
existing assumptions. The new \verb|hlPaperCorrect| predicate adds recursive
exact dependency checking without changing the faithful checker.
\verb|hl_proposition_7| and \verb|hlPaperCorrect_dependencies_unique| prove
reconstruction and uniqueness for this exact-type predicate.
\verb|hlVerifiedCorrect_has_no_dependency_reconstructor| proves that no such
reconstructor can work for all merely verified legacy proofs: two accepted
proofs can have the same stripped lines but different LEM dependencies.

\begin{remark}
The pedagogical function of the dependency column is thus not to record
information the reader could not otherwise obtain, but to require the student
to compute it. This is a defensible reason for a notation to carry redundant
data, and it is worth distinguishing from the reason a notation might carry
data that is genuinely needed.
\end{remark}

\section{Two obstructions in the other direction}\label{sec:obstructions}

\begin{definition}
A translation from Lemmon to Fitch is \emph{positional} if it preserves the
sequence of lines: the $n$th line of the Lemmon proof becomes the $n$th line
of the Fitch proof, with the same formula and the corresponding rule.
\end{definition}

Positional translation is what one wants, when it is available. It preserves
line numbers, so a student may read the two proofs side by side; and it
certifies that the two proofs are the same proof in two notations, rather than
two proofs of the same thing.

\begin{theorem}\label{thm:notpositional}
There are correct Lemmon proofs with no positional Fitch image.
\end{theorem}

We give two, illustrating two independent obstructions. The implementation
distinguishes them, and reports which was met and where:

\begin{lstlisting}
data TranslationError
  = NotNested    Int Int [Int]   -- line, assumption discharged, boxes open
  | OutOfScope   Int Int Int     -- line, line cited, box that closed over it
  | PremiseInBox Int Int         -- line, the box it was written inside
  -- ...
\end{lstlisting}

\begin{example}[Discharge order]\label{ex:order}
\begin{lemmonproof}
\dep{1}   & (1) & $P$                        & A \\
\dep{2}   & (2) & $Q$                        & A \\
\dep{1,2} & (3) & $P \wedge Q$               & 1,2 $\wedge$I \\
\dep{2}   & (4) & $P \to (P \wedge Q)$       & 1,3 CP \\
          & (5) & $Q \to (P \to (P \wedge Q))$ & 2,4 CP \\
\end{lemmonproof}
\end{example}

\noindent Assumption (1) is made first and discharged first. In a positional
Fitch image, the subproof opened by (1) would have to contain the subproof
opened by (2), since (1) is opened first and subproofs are properly nested;
but a subproof may be closed only when every subproof it contains has been
closed, and (4) closes (1) while (2) remains open. Fitch permits discharge
only in last-in-first-out order, and Lemmon imposes no such restriction.

\begin{example}[Positional trapping]\label{ex:trap}
\begin{lemmonproof}
\dep{1} & (1) & $P$                                     & A \\
\dep{2} & (2) & $Q$                                     & A \\
\dep{1} & (3) & $P \vee R$                              & 1 $\vee$I \\
\dep{2} & (4) & $Q \wedge Q$                            & 2,2 $\wedge$I \\
        & (5) & $Q \to (Q \wedge Q)$                    & 2,4 CP \\
\dep{1} & (6) & $(P \vee R) \wedge (Q \to (Q \wedge Q))$ & 3,5 $\wedge$I \\
\end{lemmonproof}
\end{example}

\noindent Here the discharge order is unobjectionable. The difficulty is that
line (3) is written between assumption (2) and its discharge at (5), and so in
any positional image lies inside the subproof opened by (2). It does not
depend on (2)---its dependency set is $\{1\}$---but that is of no help: when
the subproof closes at (5), line (3) passes out of scope, and line (6) cites
it. In the Lemmon proof nothing goes out of scope at all. Discharging an
assumption changes the dependency set of the new line; it does not affect the
availability of any earlier line.

Example~\ref{ex:trap} is the more instructive of the two, and it isolates the
second of our opening observations. The Lemmon proof does not say that line
(3) is ``inside'' anything. It is written where it is written, and that is a
fact about the page, not about the derivation. A positional translation is
obliged to read the page as though position meant something, and here it does
not.

\begin{remark}
Both examples can be repaired by permuting lines. Swapping (1) and (2) in
Example~\ref{ex:order} yields a proof in which the discharges nest; moving (3)
before (2) in Example~\ref{ex:trap} lifts it out of the subproof. This local
success does not generalise: Theorem~\ref{conj:reorder} below gives a
permutation-invariant obstruction.
\end{remark}

\section{Sharing, and two kinds of graph}\label{sec:trees}

Both obstructions concern the linear order of lines, and it is natural to ask
what becomes of them when the linear order is set aside. Doing so requires a
little graph-theoretic vocabulary, which we give in full, since it is not
universally familiar.

\subsection{Definitions}

\begin{definition}
A \emph{directed graph} consists of a set of \emph{nodes} together with a set
of ordered pairs of nodes, the \emph{edges}. We write $m \to n$ for an edge
from $m$ to $n$, and call $m$ a \emph{parent} of $n$. A \emph{path} is a
sequence of nodes each joined to the next by an edge; a \emph{cycle} is a path
returning to its starting node.
\end{definition}

\begin{definition}
A \emph{directed acyclic graph}, or \emph{DAG}, is a directed graph with no
cycles. A \emph{rooted tree} is a DAG with a distinguished node, the
\emph{root}, such that every node other than the root has exactly one parent
and every node is reachable from the root.
\end{definition}

The difference that matters here is the one in the parent count. In a tree,
each node hangs from exactly one place. In a DAG, a node may have several
parents---and when it does, we say the node is \emph{shared}: it does duty in
more than one context at once, without being written down more than once.

\begin{definition}
Let $G$ be a DAG with root $r$. Its \emph{unfolding} is the tree whose nodes
are the paths from $r$, with the obvious edges. Unfolding replaces each shared
node by one copy per context in which it is used. The unfolding of a tree is
itself; the unfolding of a DAG with genuine sharing is strictly larger, and
may be exponentially so.
\end{definition}

\subsection{What each notation presents}

The object that both notations record is a \emph{derivation}: a tree whose
root is the conclusion, whose internal nodes are applications of rules, and
whose children are the premises those rules require. This is Gentzen's
natural deduction in its original tree form, and neither Lemmon nor Fitch
writes it down as a tree. Both write it as a DAG, because both write each line
once and then refer back to it, which is precisely what sharing is.

\begin{definition}
The \emph{citation graph} of a Lemmon proof has the lines as nodes, with an
edge $m \to n$ whenever line $n$ cites line $m$.
\end{definition}

\begin{proposition}\label{prop:lemmondag}
The citation graph of a Lemmon proof is a DAG, and every line is available to
be cited by every later line.
\end{proposition}

\begin{proof}
Acyclicity is immediate, since a line may cite only lines with smaller
numbers. That every earlier line is available is a matter of the notation
containing no device by which a line could cease to be. Discharging an
assumption alters the dependency set recorded on a \emph{new} line; it does
nothing to any line already written.
\end{proof}

Fitch does have such a device.

\begin{definition}
The \emph{scope tree} of a Fitch proof has the subproofs as nodes, together
with a root representing the whole proof, and an edge from each subproof to
those immediately contained in it. A line \emph{belongs} to the innermost
subproof containing it. A line $m$ is \emph{visible} to a later line $n$ when
the node $m$ belongs to lies on the path from the root to the node $n$ belongs
to.
\end{definition}

\begin{proposition}\label{prop:fitchtree}
The scope tree of a Fitch proof is a rooted tree, and the citation graph of a
Fitch proof is a DAG in which every edge $m \to n$ has $m$ visible to $n$.
\end{proposition}

We can now say exactly where the two notations part company, and it is not
where the slogan ``Lemmon is a graph, Fitch is a tree'' would put it. Both
present a derivation as a DAG; both share. The difference is that Fitch
constrains its sharing and Lemmon does not. In a Lemmon proof any line may be
shared with any later line whatever. In a Fitch proof a line may be shared
only along the paths its scope tree permits, and once a subproof closes,
everything belonging to it is permanently unavailable.

\subsection{The translation, restated}

The two obstructions of Section~\ref{sec:obstructions} can now be seen for
what they are. By Proposition~\ref{prop:fitchtree}, translating a Lemmon proof
to Fitch means finding a scope tree under which every citation the source
makes is permitted. Both examples are
cases where no scope tree does: in Example~\ref{ex:order} because the
discharges demand incompatible nestings, in Example~\ref{ex:trap} because a
line is cited after the only subproof it could have belonged to has closed.

When no scope tree suffices, one may unfold. Unfolding removes sharing, and a
DAG without sharing is a tree; a tree imposes no constraint that could be
violated, since nothing is shared to be shared wrongly.

A derivation is represented directly as a tree. Where the Lemmon
justification names line numbers, the derivation names subderivations, and a
line cited twice becomes two subtrees.

\begin{lstlisting}
data Deriv = Deriv { dForm :: PredFormula, dRule :: DRule }

data DRule
  = DAssume Int                      -- discharged by an ancestor
  | DPremise Int                     -- never discharged
  | DMP Deriv Deriv
  | DCP Int PredFormula Deriv        -- discharges the named assumption
  | DOrE Deriv Int PredFormula Deriv Int PredFormula Deriv
  -- ...
\end{lstlisting}

Translation is then two functions and a choice between them. The positional
translation is attempted first, because it preserves numbering and structure;
unfolding is the fallback, and the result records which was used.

\begin{lstlisting}
data Route = Direct | ViaTree

lemmonToFitch :: Proof -> Either TranslationError (Route, FitchProof)
lemmonToFitch prf =
  case lemmonToFitchDirect prf of
    Right fp -> Right (Direct, fp)
    Left _   -> do d <- toDerivation prf
                   pure (ViaTree, derivationToFitch d)
\end{lstlisting}

\begin{definition}[The unfolding translation]\label{def:construction}
Given a correct Lemmon proof, unfold its citation graph into a derivation
rooted at the last line. Then traverse the derivation, emitting first the
premises, at the outermost level and once each however often they are used,
and thereafter, for each node, the subderivations its rule requires followed
by the line applying it, opening a subproof at each discharge.
\end{definition}

Here unfolding follows a line backwards to the lines it cites. Thus it uses
the reverse of the cited-to-citing edge direction in the citation-graph
definition; with that original direction the last conclusion would be a sink.

\begin{theorem}\label{thm:refuted}
The construction of Definition~\ref{def:construction} does not always yield a
correct Fitch proof.
\end{theorem}

\begin{proof}
Consider

\begin{lemmonproof}
\dep{1}   & (1) & $\forall x\, G(x)$ & A \\
\dep{2}   & (2) & $F(a)$             & A \\
\dep{1}   & (3) & $G(a)$             & 1 $\forall$E \\
\dep{1}   & (4) & $\forall y\, G(y)$ & 3 $\forall$I \\
\dep{1,2} & (5) & $F(a) \wedge \forall y\, G(y)$ & 2,4 $\wedge$I \\
\dep{1}   & (6) & $F(a) \to (F(a) \wedge \forall y\, G(y))$ & 2,5 CP \\
\end{lemmonproof}

\noindent which is correct. Line (4) generalises on $a$, and $\Gamma(4) =
\{1\}$; the sole assumption in that set has formula $\forall x\, G(x)$, in
which $a$ does not occur. The Lemmon side condition is satisfied.

The construction returns

\begin{fitchproof}
1 &                 $\forall x\, G(x)$ & Premise \\
2 & \fbar           $F(a)$             & Assume \\
3 & \fbar           $G(a)$             & $\forall$E 1 \\
4 & \fbar           $\forall y\, G(y)$ & $\forall$I 3 \\
5 & \fbar           $F(a) \wedge \forall y\, G(y)$ & $\wedge$I 2,4 \\
6 &                 $F(a) \to (F(a) \wedge \forall y\, G(y))$ & CP 2--5 \\
\end{fitchproof}

\noindent in which line 4 applies universal introduction on $a$ while $F(a)$
stands as an undischarged assumption in its scope. That is contrary to the
side condition, so the proof is not correct.
\end{proof}

It is worth being clear that this is not an artefact of the Fitch system we
chose in Section~\ref{sec:systems}. One might think to repair matters by
restating Fitch's side condition against the assumptions a line
\emph{depends on} rather than those in scope. But a Fitch line does not
record what it depends on; that is precisely the difference between the two
notations, and the reason the dependency column of a Lemmon proof, though
redundant in the sense of Proposition~\ref{prop:redundant}, is not idle.
Fitch must state its eigenvariable condition against scope, because scope is
the only thing it has. The failure is therefore a genuine one: at this rule,
Fitch's coarser bookkeeping is not merely less informative but licenses
strictly less.

Nothing is lost in the way of provability. The sequent above is Fitch-provable
by other means, and the two systems remain equivalent in the sense of proving
the same sequents. What fails is the structure-preserving translation.

There are two ways to repair it. One is to move the offending line: emit each
line at the depth of the innermost assumption in its dependency set, so that
scope and dependency coincide, and the two side conditions with them. That
proposal is not sufficient: Example~\ref{ex:trap} has disjoint, hence laminar,
dependency regions but still has no Fitch image.

The other is to move nothing and rename instead. It is available because the
Lemmon side condition says precisely what is needed to license a substitution.

\begin{lemma}[Renaming]\label{lem:rename}
Let $D$ be a derivation of $\Gamma \vdash \psi$, and let $b$ be a name
occurring nowhere in $D$. Then $D[b/a]$ is a derivation of
$\Gamma[b/a] \vdash \psi[b/a]$.
\end{lemma}

\begin{corollary}\label{cor:rename}
If in addition $a$ does not occur in $\Gamma$, then $\Gamma[b/a] = \Gamma$,
and so $\Gamma \vdash \psi[b/a]$.
\end{corollary}

Lemma~\ref{lem:rename} is an instance of the substitution theorem
of \cite[Thm.~4.5.14]{halvorson2019}: for a reconstrual $*$ with admissibility
conditions $\Delta$ for its function symbols, if $\phi \vdash \psi$ then
$\Delta, \phi^* \vdash \psi^*$. Renaming $a$ to $b$ is the reconstrual sending
$a = y$ to $b = y$, whose admissibility condition $\exists ! y\,(b = y)$ is a
theorem and so discharges, leaving $\phi^* \vdash \psi^*$.

That proof does include the case for universal introduction. It is given,
however, for a system in which generalisation is on a free variable, the side
condition being that $x$ is not free in $\Gamma$; the step goes through
because reconstrual preserves free variables, so that $x$ is not free in
$\Gamma^*$ either.

The system of \cite{halvorson2020} generalises on a name rather than a
variable, and there the step is not free. Substituting one name for another
\emph{can} introduce into $\Gamma$ the very name a universal introduction is
required to avoid---which is precisely what the free-variable lemma rules out
in the other formulation. The condition that $b$ be fresh is what closes that
gap, and since it is the step at issue here we give it.

\begin{proof}[Universal introduction]
Suppose $\Gamma \vdash \forall x\,\psi(x)$ is derived by $\forall$I from
$\Gamma \vdash \psi(c)$, where by the side condition $c$ occurs neither in
$\Gamma$ nor in $\psi(x)$, and suppose the result holds for the latter.

If $c \neq a$, then $\psi(c)[b/a] = (\psi[b/a])(c)$, and the inductive
hypothesis gives $\Gamma[b/a] \vdash (\psi[b/a])(c)$. Universal introduction
on $c$ applies: $c$ does not occur in $\Gamma[b/a]$, since it did not occur in
$\Gamma$ and the substitution introduces only $b$, which differs from $c$
because $b$ occurs nowhere in $D$ while $c$ does; and $c$ does not occur in
$\psi(x)[b/a]$ for the same reason. So
$\Gamma[b/a] \vdash \forall x\,(\psi[b/a])(x)$, which is
$(\forall x\,\psi(x))[b/a]$.

If $c = a$, then $a$ occurs neither in $\Gamma$ nor in $\psi(x)$, so
$\Gamma[b/a] = \Gamma$ and $\psi(a)[b/a] = \psi(b)$. The inductive hypothesis
gives $\Gamma \vdash \psi(b)$. Universal introduction on $b$ applies: $b$
occurs neither in $\Gamma$ nor in $\psi(x)$, since it occurs nowhere in $D$ at
all. So $\Gamma \vdash \forall x\,\psi(x)$, and since $a$ does not occur in
$\forall x\,\psi(x)$, this is $(\forall x\,\psi(x))[b/a]$.
\end{proof}

Existential elimination goes the same way, its witness condition being of the
same shape; the remaining rules impose no condition on names and follow the
pattern of the propositional cases.

\begin{remark}\label{rem:freshness}
The freshness of $b$ is not a convenience, and the contrast with the
free-variable formulation shows why. There, reconstrual preserves free
variables, so a universal introduction licensed before the substitution is
licensed after it. Here, substitution of one name for another can \emph{add} a
name to $\Gamma$: if $b$ already occurs in $\Gamma$, a universal introduction
on $b$ somewhere within $D$ may be licensed before the substitution and not
after, since the substitution places $b$ among the assumptions that
introduction is required to avoid. It is the second case of the proof above
that this breaks.
\end{remark}

The repair follows, by Corollary~\ref{cor:rename}. Wherever the construction
is about to emit a universal introduction whose eigenvariable occurs in an
assumption in scope, rename the subderivation feeding it to a name occurring
nowhere in the proof---nowhere, and not merely nowhere in the subderivation,
for the reason given in Remark~\ref{rem:freshness}. Its open assumptions are
untouched, since by the Lemmon side condition $a$ does not occur in them. Its conclusion is untouched, since abstraction removes every
occurrence of $a$. And the new name occurs in no assumption whatever, so
\emph{a fortiori} in none in scope. Note that the unfolding shares nothing, so
renaming within one subderivation cannot disturb any other.

Applied to the proof of Theorem~\ref{thm:refuted}, a single line changes:

\begin{fitchproof}
1 &                 $\forall x\, G(x)$ & Premise \\
2 & \fbar           $F(a)$             & Assume \\
3 & \fbar           $G(b)$             & $\forall$E 1 \\
4 & \fbar           $\forall y\, G(y)$ & $\forall$I 3 \\
5 & \fbar           $F(a) \wedge \forall y\, G(y)$ & $\wedge$I 2,4 \\
6 &                 $F(a) \to (F(a) \wedge \forall y\, G(y))$ & CP 2--5 \\
\end{fitchproof}

\noindent Line 3 instantiates to $b$ rather than $a$, and line 4 generalises
on a name that occurs in no assumption anywhere. Nothing else moves.

\begin{theorem}[Totality of the repaired unfolding]\label{conj:total}
For the compact reconstructed calculus, modify the construction of
Definition~\ref{def:construction} by globally fresh renaming at both universal
introduction and existential elimination whenever the relevant name occurs in
an assumption in scope. This construction yields for every nonempty correct
Lemmon proof a well-formed and correct Fitch proof with the same conclusion.
\end{theorem}

This is a corrected replacement for original Conjecture 26. That conjecture
explicitly concerned Definition 21's algorithm with universal-introduction
repair alone; it was not merely a statement that the systems prove the same
sequents. The existential-elimination counterexample requires the additional
repair. The sequent-containment theorem in Isabelle is a consequence of the
corrected construction theorem.

Two things recommend this repair over the first. It is local within each
unfolded subderivation: no additional rearrangement or change of subproof shape
is required. And it makes no assumption that a duplication-free image exists.
The Isabelle proof packages the four structural lemmas of
Section~\ref{sec:open} with the verified renaming step.

The construction has a cost, which is the cost of unfolding: a line shared $k$
times is derived $k$ times, and in the worst case the result is exponential in
the size of the source. There is also a small point of bookkeeping. If a
subproof's conclusion is a line from outside it---as in deriving $Q \to P$
from the premise $P$---then the subproof would contain nothing at all, and
Fitch requires a reiteration step to bring the line in. Lemmon has no such
rule and needs none, by Proposition~\ref{prop:lemmondag}; and the reiterated
line is a tautological consequence of the line it repeats, so nothing is added
to the system.

\section{Formal answers}\label{sec:open}

\subsection{The totality theorem}

The four obligations suggested by the construction are machine-checked for
the compact calculus. In the Isabelle sources they appear as
\verb|unfolding_terminates|, \verb|L2_L3|, \verb|L4_derivation|, and
\verb|conjecture_26|:

\begin{enumerate}
\item[(L1)] \emph{Unfolding terminates.} A correct, nonempty Lemmon proof is
converted to a finite derivation tree. The recursion is well-founded because
citations point to smaller line numbers.

\item[(L2)] \emph{Citations precede uses.} The traversal emits every
subderivation before the rule application that uses it.

\item[(L3)] \emph{Every citation is in scope.} Premises are emitted at the
outermost level, while a discharged assumption heads the subproof containing
the uses discharged with it. Isabelle proves the complete resulting Fitch
object structurally well formed.

\item[(L4)] \emph{Each rule transfers.} The generated line satisfies the
corresponding Fitch rule. The proof includes the subproof endpoints required
by conditional proof, reductio, disjunction elimination and existential
elimination. It also checks the freshness repair for universal introduction
and the analogous witness condition for existential elimination.
\end{enumerate}

The final theorem is stronger than termination: it gives a derivation $D$,
proves the generated Fitch proof correct, and proves that its conclusion is
the conclusion of the source Lemmon proof. This is
Theorem~\ref{conj:total}, the corrected compact replacement for the original
totality conjecture.

For the exact \verb|HL_| language, \verb|hlUnfoldDerivation_succeeds| and
\verb|hlToDerivation_total| now establish L1, including sufficient recursion
fuel and success on every correct nonempty verified source.
\verb|hlToDerivation_conclusion| preserves its selected conclusion.
The general counterparts of L2--L4, syntactic preservation by fresh renaming,
and correctness of the unfolded derivation remain to be proved.

The theorems \verb|hlViaTree_accepted| and
\verb|hlLemmonToFitchChecked_accepted| establish Fitch correctness,
conclusion preservation and premise inclusion whenever
the checked operation returns a proof. An output guard can ensure that
every returned proof is correct; it cannot establish that every correct source
produces an accepted target. Nor does the semantic soundness theorem for all
twenty-one Lemmon rules prove that their emitted Fitch applications are correct.

\subsection{Why the eigenvariable rules need attention}\label{sec:eigen}

Theorem~\ref{thm:refuted} turns on universal introduction, but the same audit
must be made for existential elimination: its witness name also carries a
freshness condition involving assumptions. The remaining rules place their
substantive conditions on the formulas cited and the formula concluded.
Arbitrariness, by contrast, is a matter of what a line rests on. Lemmon can
say this exactly, because a Lemmon line records its dependencies. Fitch can
only say it about scope. Wherever scope properly includes the dependency
set---which by Proposition~\ref{prop:minimal} it may---the two conditions can
come apart. The formal rule-transfer proof checks both eigenvariable cases.

That the repair is a renaming rather than a rearrangement is a point in
favour of the diagnosis. The difficulty is not that the line is in the wrong
place; it is that the \emph{name} is doing double duty, standing both as the
arbitrary individual of the generalisation and as the individual the enclosing
assumption is about. Lemmon can tell those two roles apart because it knows
which assumptions the line answers to. Fitch cannot, and so a name that is
arbitrary in the one sense is not visibly arbitrary in the other. Giving the
generalisation a name of its own removes the ambiguity without touching the
structure of the proof at all.

\subsection{Negative answers to the economy questions}

\begin{theorem}[Permutation does not always suffice]\label{conj:reorder}
There is a correct Lemmon proof such that no citation-respecting permutation
of its lines has a positional Fitch image.
\end{theorem}

\begin{theorem}[Duplication-free images do not always exist]
\label{conj:noduplication}
There is a correct Lemmon proof with no well-formed Fitch image containing
each source line exactly once while preserving its formula and corresponding
citation-bearing rule.
\end{theorem}

Both negative answers have the same five-line witness:

\begin{lemmonproof}
\dep{1}   & (1) & $P$                         & A \\
\dep{2}   & (2) & $Q$                         & A \\
\dep{1,2} & (3) & $P \wedge Q$                & 1,2 $\wedge$I \\
\dep{2}   & (4) & $P \to (P \wedge Q)$        & 1,3 CP \\
\dep{1}   & (5) & $Q \to (P \wedge Q)$        & 2,3 CP \\
\end{lemmonproof}

\noindent Lines (4) and (5) demand subproof references $(1,3)$ and $(2,3)$:
two different assumptions are discharged at the same conclusion line. But in
a well-formed Fitch proof two distinct subproofs cannot have the same last
line. This obstruction is invariant under permutation, proving
Theorem~\ref{conj:reorder}; and it prevents any arrangement containing each
source line once, proving Theorem~\ref{conj:noduplication}.

The Isabelle result is stronger. Let a target \emph{cover} the source when
every source line occurs with its number, formula, and corresponding
citation-bearing rule, while arbitrary additional target lines are allowed.
Even then the displayed proof has no well-formed Fitch image. Auxiliary lines
cannot change the fact that the two source discharge rules require different
subproofs ending at line (3). Thus duplication or administrative additions
alone cannot repair the obstruction; a more liberal translation must unfold
or otherwise rewrite source occurrences.

Original Conjecture 28 did not formally define a nonpositional image. The
negative theorem uses the correspondence just stated: source discharge rules
retain their prescribed endpoints, even when auxiliary lines are allowed.
It does not settle a broader translation relation that redirects those
citations to auxiliary derivations. The Conjecture 27 theorem, by contrast,
explicitly quantifies over all injective renumberings, including every
citation-respecting permutation.

There is a narrower negative result about laminarity. The dependency regions
of Example~\ref{ex:trap} are disjoint, hence laminar, but
\verb|theorem_10_ex12| rules out a Fitch image preserving the original
numbered lines. Strictly increasing target numbering retains the original
order. This proves laminarity insufficient for that fixed-order image; it
does not refute sufficiency after renumbering and reordering. Indeed the
original Remark 13 explicitly repairs this example by permutation.

\section{Machine verification}\label{sec:machine}

The development is formalised in Isabelle/HOL. Formulas, proof lines, Fitch
subproofs, all rule checks, both translations, and the predicates used in the
negative results are executable definitions. For example, the paper's word
\emph{positional} is represented by an ordinary predicate, not by an informal
label:

\begin{lstlisting}[language={}]
definition positionalImage ::
  "fitch_proof => lemmon_proof => bool" where
  "positionalImage F P =
     list_all2 lineMatches (flatten F) P"
\end{lstlisting}

\noindent Here \verb|lineMatches| requires equality of line number and formula
and the corresponding rule; \verb|fitchWF F| is asserted separately whenever
the target is required to be an actual Fitch proof. Keeping those two clauses
separate proved useful: without well-formedness, any list of Lemmon lines could
be laid out flat and called a positional image.

The repaired totality claim is a theorem over arbitrary nonempty correct
source proofs of the compact calculus:

\begin{lstlisting}[language={}]
theorem conjecture_26:
  assumes "lemmonCorrect P" and "P \<noteq> []"
  shows "\<exists>d. toDerivation P = Inr d
       \<and> fitchCorrect (derivationToFitch d)
       \<and> fitchConclusion (derivationToFitch d)
           = conclusion P"
\end{lstlisting}

\noindent The retained name \verb|conjecture_26| records its correspondence
with the earlier draft; its Isabelle status is \verb|theorem|, and its
construction includes both eigenconstant repairs. The other compact
machine-checked answers are recorded under the following source names:

\begin{lstlisting}[language={}]
theorem_10
conjecture_27_false
conjecture_28_with_auxiliaries_false
\end{lstlisting}

\noindent The first proves the two original witnesses have no positional
image, the second refutes the permutation claim, and the third proves the
stronger auxiliary-line result. This is one reason the formal work is helpful
philosophically: apparently modest changes in the meaning of \emph{same
proof}---same order, same occurrences, or merely coverage of source
lines---become distinct mathematical claims.

The direct translation reports seven distinct obstructions, and it is natural
to ask whether they are necessary as well as sufficient: does a source they
reject really have no positional image? For every witness the checks were
built from the answer is yes. The combined theorem is
\begin{lstlisting}[language={}]
checks_reject_only_imageless_sources
\end{lstlisting}
\noindent The six sources exercise nesting, line visibility, subproof
visibility, premise order, discharge order and assumption reuse. Each is
proved to have no well-formed positional Fitch image.
Two of those proofs are short and turn on facts already available: a cited
subproof runs forwards, because its assumption line lies in its own span
(\verb|image_subproof_ordered|); and a subproof is determined by its
assumption line (\verb|subs_unique_proof|), so one assumption cannot head two
subproofs closing at different lines.

This is evidence for the converse, not a proof of it: six sources are not all
sources. An earlier counterexample to the converse did exist, representing a
late undischarged assumption as an unused one-line subproof; it ended inside
that subproof, and the closed-root condition of the previous subsection
removes it. We record the question as open.

\subsection{The authoritative \emph{How Logic Works} checker}

The Haskell source is not treated as disposable prototype code. Its logical
datatypes and rule decisions are the specification for the executable
\emph{How Logic Works} checker. Isabelle represents them without first
translating them into the compact formula and rule types used in the earlier
proofs:

\begin{lstlisting}[language={}]
datatype hl_formula =
    HL_Predicate hl_name "hl_term list"
  | HL_Boolean bool
  | HL_Not hl_formula
  | HL_And hl_formula hl_formula
  | HL_Or hl_formula hl_formula
  | HL_Implies hl_formula hl_formula
  | HL_Iff hl_formula hl_formula
  | HL_ForAll hl_name hl_formula
  | HL_Exists hl_name hl_formula
\end{lstlisting}

\noindent The corresponding \verb|hl_justification| datatype has exactly the
twenty-one Haskell constructors. Line numbers and dependency sets use
integers, as Haskell's \verb|Int| does. The predicate \verb|hlCorrect|
reproduces the source checker's full-proof lookup, uniqueness, forward-number
test, formula patterns, and dependency equations. A strengthened
\verb|hlVerifiedCorrect| additionally requires positive increasing source
order and that every dependency really names an assumption. This distinction
is deliberate: exact reproduction makes source-level comparison possible,
while the strengthened predicate states the invariant needed for theorems.

Both added invariants are genuinely extra, and each has a witness accepted by
\verb|hlCorrect| and rejected by \verb|hlVerifiedCorrect|.
\verb|hl_cp_misordered| is a worked example with two of its lines exchanged in
the list and their numbers left alone: source order is the only condition it
fails, because the supplied \verb|checkStructure| compares a cited line's
number with the citing line's number rather than checking physical order.
\verb|hl_lem_free_dependency| fails only dependency closure, because excluded
middle is the one rule of the source checker that constrains nothing about the
dependency set, so such a line may name a line that is not an assumption at
all. Neither witness is unsound---each concludes a tautology---and we have not
settled whether \verb|hlCorrect| alone admits a semantically invalid proof.
What can be said is why the obvious attacks fail: every rule but excluded
middle fixes the dependency set from the cited lines; excluded middle leaves
it free but proves a validity, so a free set can only enlarge the reported
premises and thereby weaken the sequent; and every subtraction discharges a
line the checker separately requires to be an assumption.

The Haskell-shaped Fitch datatype is likewise represented directly. Its total
map to Lemmon notation is written with the same delta:

\begin{lstlisting}[language={}]
abbreviation hlDelta ::
  "hl_fitch_proof => hl_proof" ("\<delta>\<^sub>H _") where
  "\<delta>\<^sub>H F == hlFitchToLemmon F"
\end{lstlisting}

\noindent Reiteration becomes \verb|HL_PropTaut [m]| under this map, just as
in \verb|FitchConvert.hs|. Executable Isabelle lemmas check representative
modus-ponens and conditional-proof images and one accepted proof for every
authoritative justification constructor.

The converse direct route from \verb|FitchConvert.hs| is now represented on
the same source-shaped datatypes as \verb|hlLemmonToFitchDirect|.  A successful
result is proved positional: flattening it preserves every source line number
and formula in order.  The Isabelle version also incorporates the structural
repairs established by the compact development: reversed or multiply reused
boxes, late premises, bad subproof scope, and the stronger Fitch
eigenvariable-in-scope condition are diagnosed before an image is returned.
The executable conditional-proof example verifies both the generated Fitch
proof and the round trip under $\delta_H$.

The derivation-tree route from \verb|FitchConvert.hs| is represented on these
same datatypes as well.  The public operation
\verb|hlLemmonToFitchChecked| accepts a direct image only when
\verb|hlFitchCorrect| verifies it, its conclusion matches the source, and its
displayed premises form a subset of the source conclusion's open premises.
Otherwise it unfolds, repairs eigenconstants and lays out nested subproofs;
the fallback must pass those same three checks. In particular, unused source
assumptions cannot silently enlarge the translated sequent.
The repair scope includes all emitted outer premises, as well as enclosing
box assumptions. This fixes the case in which a conflicting constant occurs
in an outer premise rather than a box. Both eigenconstant rules are checked.
An executable shared-discharge example forces the fallback. The checked
operation may report failure: successful unfolding alone does not prove that
every candidate passes these guards. The raw direct/diagnostic interfaces do
not carry the checked operation's full guarantee.

That example exposed one further defect in the handwritten fallback.  Its
single global \verb|boxed| set labels an assumption as discharged if any line
of the source proof discharges it, even when that line lies on a dead sibling
branch omitted from the derivation of the selected conclusion.  The resulting
tree can drop a genuine open premise and later cite line zero.  Isabelle's
\verb|hlClassifyAssumptions| instead carries the dischargers on the actual
ancestor path, which is the condition the tree representation requires.

\subsection{Executable tests and the regression corpus}

The original translations were implemented in Haskell against the proof
checker for \emph{How Logic Works}. Those tests remain useful as concrete
validation of the parser, checker, renderer, and examples. They do not
establish the missing general HL construction theorem. The
Isabelle mirror now shares their formula and justification structure. The
existing application has not yet been mechanically linked against the
generated module, so application-level tests and Isabelle-level executable
lemmas remain separate regression boundaries.

The current regression source contains $80$ cases, of which $29$ are marked
valid. Twenty-four are expected to admit a positional translation, and for
each the round trip is required to return the original proof line for line,
\emph{including its dependency sets}---an instance of
Proposition~\ref{prop:redundant}, since the sets were recomputed and not
copied. Five are expected to use the derivation-tree route: the obstruction
examples, a premise placed inside a subproof, an ordinary reductio whose
premise occurs after the discharged assumption, and the eigenvariable example.
For those cases the suite requires the recovered Lemmon proof to check and to
have the same conclusion, rather than to preserve line numbers.

Sixteen further proofs, transcribed from photographs of student work, were
translated. Fourteen were well formed; all fourteen admitted positional
translations and round-tripped exactly. The remaining two contained malformed
justifications on the page---a rule cited with the wrong number of
lines---and were rejected by the parser, as they should be.

We record the empirical finding for what it is worth: neither obstruction
occurred in any student proof. They are not artefacts, as
Section~\ref{sec:obstructions} shows, but they appear to be rare in practice.

One methodological point deserves recording, because it cost us two false
negatives. The round trip is a weaker test than it appears. Since $\delta$
recomputes dependency sets from the rules rather than reading them off the
page, a Fitch proof that is \emph{malformed}---one violating the conditions in
Definition~\ref{def:fitch}---may still return the Lemmon proof it came from,
exactly. Two proofs in the corpus did precisely this, both by placing a
premise inside a subproof, where its dependency set recomputes correctly while
the subproof appears to derive the premise from its assumption. Neither was
detected by the round trip; the first was noticed by looking at the rendered
output. The remedy is to check the conditions of the notation directly, and
separately:

\begin{lstlisting}
-- Conditions the notation imposes, which a round trip cannot detect:
--   * a premise may occur only at the outermost level;
--   * a line may cite only what is in scope.
fitchWellFormed :: FitchProof -> Maybe String
\end{lstlisting}

\noindent The moral generalises beyond this case. A test that verifies an
invariant is preserved under translation cannot, by itself, establish that
either side satisfies the conditions of its own notation.

\subsection{Semantic audit}

The Isabelle development gives both formal layers a standard model-theoretic
semantics.  For the compact internal language it proves soundness of every
derivation-tree rule (\texttt{derivOK\allowbreak\_sound}) and semantic soundness of the
Lemmon turnstile (\verb|lemmon_semantic_soundness|).  More importantly for the
executable checker, the Haskell-shaped layer now has a complete soundness
argument of its own.  The local theorem \verb|hlRuleOK_sound| covers every one
of the twenty-one authoritative justification constructors.  A well-founded
induction over cited line numbers then proves
\texttt{hlVerifiedCorrect\allowbreak\_line\allowbreak\_sound}, and the
resulting end-to-end theorem
\texttt{hlVerifiedCorrect\allowbreak\_sound} says that
if all open premises are true, the checked conclusion is true in every standard
interpretation.

The induction is stronger than a fixed-model argument.  Universal
introduction and existential elimination temporarily reinterpret a fresh
constant, so the induction hypothesis must apply uniformly to every standard
interpretation on the same domain.  Making that requirement explicit was one
of the useful consequences of formalising the checker rather than merely
testing representative proofs.

The Haskell-shaped semantics must also treat the predicate name ``='' as
identity for equality introduction and elimination to be sound; the
handwritten \verb|ModelSemantics.hs| currently treats it as an arbitrary
predicate relation. Formalising the authoritative layer makes that hidden
model condition explicit. The Isabelle theorem suite records the principal
checks under the following names:

\begin{lstlisting}[language={}]
hlPropositionalConsequence_sound
hl_haskell_model_semantics_needs_equality_condition
hlQuantifierNegationEquivalent_sound
hlEqualityElimination_sound
hlRuleOK_sound
hlVerifiedCorrect_line_sound
hlVerifiedCorrect_sound
\end{lstlisting}

\noindent The first connects the executable \verb|PropTaut| valuation search
to arbitrary interpretations. The second puts the accepted equality-
introduction proof beside a refuting arbitrary-equality model. The last three
lines of the list give local, line-by-line, and complete-proof soundness,
respectively; the intervening specialized lemmas discharge quantifier
negation and equality replacement inside the local proof.

For executable code Isabelle keeps this mathematical interpretation separate
from the finite, error-aware representation

\begin{lstlisting}[language={}]
HL_FiniteModel
\end{lstlisting}

\noindent corresponding to \verb|ModelSemantics.hs|. Its evaluator returns
\verb|None| for an unbound
variable, a missing interpretation, or an empty quantifier domain, and it
handles biconditionals explicitly.  The companion \verb|hlToDNF| operation
reproduces the propositional DNF pipeline while repairing the omitted
biconditional traversal.

The same audit found a second, independent boundary case at the root of a
Fitch proof. The one-item proof consisting of a single empty subproof satisfied
every well-formedness condition then imposed, reported no premises, and
reported that subproof's assumption as its conclusion. The turnstile therefore
derived an arbitrary $P$ from nothing, and any interpretation making $P$ false
refuted the sequent. Nothing among the conditions looked at where the proof
ends, and an empty subproof body satisfies ``a subproof ends in a line''
vacuously.

The new condition, \verb|concludesAtTop|, requires the last item of a proof
to be a line at the
outermost level. It is not redundant:
\verb|openAssumptionFitch_fails_only_root_scope| records by evaluation that the
witness satisfies each of the other six conditions and fails only this one.
With the conjunct in force the Fitch turnstile is semantically sound
(\verb|fitch_semantic_soundness|), and the argument needs no new model theory:
\verb|theorem_4_forward| carries the proof to Lemmon, \verb|proposition_5|
bounds the conclusion's dependencies by the assumptions in scope at it, and the
new conjunct makes that scope the premises alone.

The supplied Haskell has the same omission: \verb|fitchWellFormed| checks scope
and citations but never asks where the proof ends. There the exact layer keeps
\verb|hlFitchWellFormed| faithful to the source and adds the condition
separately, as \verb|hlConcludesAtTop|, which the verified predicates
\verb|hlFitchVerified| and \verb|hlFitchCorrect| require. Formalisation is
valuable here for exposing a convention that both the informal presentation and
the running code had left implicit, and then for saying exactly what supplying
it costs and what it buys.

A further review found a different scope failure even with a final top-level
line: the faithful structural checker exposes a closed subproof's assumption
and last line as ordinary lines outside it. The strengthened predicates now
add \verb|hlFitchNestedWellFormed|, which checks actual ordinary-line visibility
and genuine subproof references independently. Eigenconstant scope is recorded
from outer premises and enclosing assumptions, without adding closed boxes or
earlier derived formulas. A separate \verb|hlFitchPremiseClosed| guard connects
the erasure's final dependencies to the displayed outer premises.
\verb|hlFitchVerified_sound| and \verb|hlFitchCorrect_sound| then prove the
reported Fitch conclusion from those premises in every standard interpretation.
This proof uses the explicit premise-closure guard; it does not prove that
the guard follows from nesting alone or that the emitter always satisfies it.

\subsection{Generated Haskell and SML}

The cap theory \verb|Lemmon_Fitch.thy| is also the code-generation entry point.
Isabelle exports Haskell and SML from the same executable constants used in
the proofs. The generated API now contains two clearly named layers: the
compact translation development, and the authoritative \verb|HL_| datatypes,
checker, finite evaluator, DNF conversion, $\delta_H$, repaired direct
translation, and complete derivation-tree fallback. Generated files
are derived artefacts; the Haskell program is authoritative for the intended
\emph{How Logic Works} rule choices, while Isabelle is the authoritative,
total specification and proof environment for their formal reproduction.

The source audit also found partial cases that a total Isabelle function could
not silently inherit. \verb|ProofTypes.hs| omitted \verb|Boolean| from several
recursive operations; \texttt{ModelSemantics.}\allowbreak\texttt{eval}
omitted biconditionals; and
\verb|PropDNF| omitted biconditionals from its traversal. Isabelle supplies
all of those cases. It also exposes two structural choices in the handwritten
checker: LEM does not constrain its dependency set, and the forward-reference
test compares numbers without requiring the file itself to be in increasing
order. The exact predicate records those behaviors; the verified predicate
adds the intended invariants.

Copying the generated module beside the web application still does not replace
its parser, JSON codecs, OCR/transcription pipeline, LaTeX rendering, web
routes, static pages, or corpus. Those are interface code, not products of
Isabelle's code generator. An application adapter is therefore still needed,
but no further rule-roster decision is pending: the supplied Haskell roster is
the one to which that adapter must connect.

The two checkers have been compared directly. Both were compiled --- the
handwritten kernel unmodified --- and run on the same inputs, with verdicts
compared line by line through the public \verb|checkProof| interface. Across
9{,}874 comparisons there was no disagreement: the fourteen proofs of the
\verb|eval| corpus that parse; twelve further proofs written to exercise the
eight rules that corpus never uses, among them excluded middle, propositional
consequence, the biconditional rules and quantifier negation; and 1{,}728
systematic mutants which corrupt dependency sets, line numbers, justifications
and formulas, so that 2{,}706 of the comparisons are rejections rather than
acceptances. The mutants are the point: a corpus of valid proofs would be
passed by a checker that accepted everything.

This is testing rather than proof, and cannot be otherwise, since Haskell has
no formal semantics here. Two further points are worth recording. Attempting
to compile against the generated module showed that it had exported its
integer and set types abstractly, so that no caller could construct an
argument to any exported checker; the omission had gone unnoticed because
nothing had previously been compiled against it. And one divergence is
permanent: the handwritten line numbers are machine integers and wrap, while
generated ones are unbounded, because Isabelle will not emit a wrapping
integer without falsifying the proofs about it.

\section{What kind of equivalence is this?}\label{sec:equivalence}

It is tempting to summarise all of this by saying that Lemmon and Fitch
notation are intertranslatable, and to leave it there. The results above
suggest that the summary conceals more than it reports.

For the compact calculus, translation from correct Fitch proofs to Lemmon
proofs is total, canonical,
and not injective (Theorem~\ref{thm:total},
Corollary~\ref{cor:notinjective}). Translation from Lemmon to Fitch is not
positional (Theorem~\ref{thm:notpositional}), but it is total by way of the
repaired unfolding (Theorem~\ref{conj:total}). Theorems
\ref{conj:reorder} and \ref{conj:noduplication} show that neither permutation
nor a target preserving each source citation-bearing rule once can serve in
general under the explicit image relation used here.

It would be premature to describe the current formal objects as a retract.
The failure of the original correctness biconditional already prevents its
stated argument. In addition, unfolding may duplicate or omit source lines,
insert reiterations and rename constants, so it does not give an identity
round trip on raw proof objects. A retract or normalization theorem would
require an explicit equivalence relation and the appropriate equations; none
has been proved here. What is already proved for the compact calculus is that
every correct Lemmon proof has a correct Fitch derivation with the same
conclusion, and every correct Fitch proof has a correct Lemmon image. For the
exact HL language, the general emitter correctness and totality theorem
remains an explicit completion requirement.

The reason is that the two notations record different amounts. A Lemmon proof
carries its dependency structure exactly, and carries no positional
information at all. A Fitch proof carries its dependency structure only up to
an upper bound, and carries positional information in addition---information
which is, from the point of view of the derivation, arbitrary. Neither is a
richer notation than the other. They are differently lossy.

That is a more interesting relationship than notational variance, and it is
the sort of relationship that arises whenever two formalisms are said to say
the same thing. The case is unusually clean: both formalisms are small,
completely specified, and mechanically checkable, so the question of what
survives translation can be settled rather than argued. We suggest it is worth
having such a case in view when the same question is asked of theories, where
it cannot.

\begin{thebibliography}{9}

\bibitem{halvorson2020}
Hans Halvorson.
\emph{How Logic Works: A User's Guide}.
Princeton University Press, 2020.

\bibitem{halvorson2019}
Hans Halvorson.
\emph{The Logic in Philosophy of Science}.
Cambridge University Press, 2019.

\bibitem{lemmon1965}
E.~J. Lemmon.
\emph{Beginning Logic}.
Nelson, 1965.

\bibitem{fitch1952}
Frederic B. Fitch.
\emph{Symbolic Logic: An Introduction}.
Ronald Press, 1952.

\bibitem{suppes1957}
Patrick Suppes.
\emph{Introduction to Logic}.
Van Nostrand, 1957.

\end{thebibliography}

\end{document}
